\textcolor{black}{\section{Conclusiones}}

\noindent Para la segunda parte del experimento se puede observar rápidamente que, en cada uno de los tres voltajes de fuente distintos, el $V_{rb}$ tiene una tendencia al aumento conforme al aumento del $V_c$, mientras que el $V_a$ presenta un tipo de máximo.

Los electrones emitidos por calentar el cátodo son acelerados por aumentarle a él mismo el voltaje, pero son dispersados por las moléculas del gas, y la mayoría son recogidas por la rejilla-blindaje. El máximo del $V_a$ representa el momento en el que la dispersión del gas fue mínima, ya que llegan al ánodo solo los electrones que lograron atravesar el gas sin ser dispersados.

Lo anterior explica el aumento del $V_{rb}$, que se nota creciente conforme al $V_c$ e incluso se observa que llega a superar al $V_a$ en las Figuras \ref{Graf:3v_manual} y \ref{Graf:4v_manual}. También es posible notar que la dispersión mínima del gas a un $V_f = 2\,V$ se da alrededor de los $2.2\,V$, mientras que a $V_f = 3\,V$ y $V_f = 4\,V$ se observa en $1.7\,V$ y $1.5\,V$, respectivamente.

Finalmente, las Figuras \ref{Graf:2v_automatic}, \ref{Graf:3v_automatic} y \ref{Graf:4v_automatic} muestran demasiada dispersión en las mediciones del $V_a$, por lo que se recomendaría modificarlo a que tenga un tiempo pequeño de espera antes de hacer la medición para que el sistema se estabilice, o también que tome varias lecturas en ese mismo $V_c$ y entregue, como lectura, el promedio de ellos.